\documentclass[10pt,a4paper]{amsart}
\usepackage[margin=19mm]{geometry}
\usepackage{amsmath,amssymb,mathtools}
\usepackage[scr=rsfs,cal=euler]{mathalfa}
\usepackage{xfrac}
\usepackage[unicode,hidelinks,bookmarks=false]{hyperref}
\newcommand{\ie}{i.e.\ }
\newcommand{\cf}{cf.\ }
\newcommand{\resp}{respectively\ }
\newcommand{\Subsection}{\subsection}
\providecommand{\nobreakdash}{\nobreak\hbox{-}}
\setlength{\emergencystretch}{2em}
\multlinegap0pt
\usepackage{natbib}
\usepackage{mathtools}
\usepackage{amssymb}
\usepackage{bm}
\usepackage{siunitx}
\usepackage{booktabs}
\usepackage{multirow}
\usepackage{caption}
\usepackage{float}
\usepackage[normalem]{ulem}
\usepackage[shortlabels]{enumitem}
\usepackage{xcolor}
\usepackage{algorithm}\usepackage{algpseudocode}
\usepackage{tikz}
\newtheorem{definition}{Definition}
\newtheorem{theorem}{Theorem}
\title{VERIFY-RL: Verifiable Recursive Decomposition for Reinforcement Learning in Mathematical Reasoning}
\author{Kaleem Ullah Qasim, Jiashu Zhang, Hao Li, Muhammad Kafeel Shaheen}
\date{Source: arXiv:2602.07559v1 (CC BY 4.0)}
\begin{document}
\maketitle

\noindent\emph{Source and licence.} VERIFY-RL: Verifiable Recursive Decomposition for Reinforcement Learning in Mathematical Reasoning. Version arXiv:2602.07559v1 (CC BY 4.0), distributed by arXiv under the Creative Commons Attribution 4.0 International licence, http://creativecommons.org/licenses/by/4.0/, as recorded in the arXiv licence metadata for this exact version and shown on the abstract page https://arxiv.org/abs/2602.07559v1. Full licence notice: \texttt{licenses/arxiv-2602.07559-CC-BY-4.0.txt}.\par\smallskip

\noindent\textit{Editorial scope.} Counted unit: the article's theorem thm:chain (source0.tex lines 425--428). The article's complete original proof of it is delivered with it (source0.tex lines 430--434, one \texttt{proof} environment of 5 lines). No mathematics is written, repaired or re-derived: the statement and the proof are the article's own lines, in the article's own notation, and no retained line is substituted or re-flowed.  Retained uncounted as the source-local context the proof needs: lines 255--265; lines 267--277.  Nothing was omitted from any retained span, so the package carries no omission record.  No retained line contains a citation, so the wrapper carries no bibliography and no bibliography entry.  No retained line contains a figure environment, an image-inclusion command or drawing code, so no image is delivered and the package declares no asset dependency.  Editorial formatting only, in addition: an AMS \texttt{amsart} preamble that restates the article's own declarations and macros used by the retained lines, namely the environments \texttt{definition}, \texttt{theorem} on separate counters, together with the standard packages those lines need.  The retained environments print in this document as Definition 1, Theorem 1 in order of appearance, whereas the article prints the same items with the numbering of its own full document.  The complete native source of the article as distributed in the arXiv e-print for this exact version is bundled unchanged as \texttt{sources/194188/source0.tex}.
\begin{definition}[Expression Tree]
\label{def:expr_tree}
For $f \in \mathcal{F}$, the expression tree $T_f$ is defined inductively:
\begin{align}
T_x &= \text{leaf}(x) \\
T_c &= \text{leaf}(c) \quad \text{for } c \in \mathbb{R} \\
T_{f \circ g} &= \text{node}(f, [T_g]) \\
T_{f \cdot g} &= \text{node}(\times, [T_f, T_g]) \\
T_{f + g} &= \text{node}(+, [T_f, T_g])
\end{align}
\end{definition}

\begin{definition}[Nesting Depth]
\label{def:depth}
The nesting depth $\delta: \mathcal{F} \rightarrow \mathbb{N}$ measures structural complexity:
\begin{equation}
\delta(e) = \begin{cases}
0 & \text{if } e \in \{x\} \cup \mathbb{R} \\[4pt]
1 + \max\limits_{i \in [n]} \delta(a_i) & \text{if } e = g(a_1, \ldots, a_n) \text{ for } g \in \mathcal{B} \\[4pt]
\max\limits_{i \in [n]} \delta(a_i) & \text{if } e = \sum_{i=1}^n a_i \text{ or } e = \prod_{i=1}^n a_i
\end{cases}
\end{equation}
\end{definition}

\begin{theorem}[Chain Rule Satisfies V1\textendash V3]
\label{thm:chain}
For $p = f(g(x))$ with $g(x) \neq x$ and $c = g(x)$: $\delta(c) < \delta(p)$, $\sigma(c) \preceq \sigma(p)$, and $c \sqsubseteq p$.
\end{theorem}

\begin{proof}
\textbf{(V1)} By Definition~\ref{def:depth}: $\delta(p) = \delta(f(g(x))) = 1 + \delta(g(x)) = 1 + \delta(c)$. Since $1 + \delta(c) > \delta(c)$, we have $\delta(p) > \delta(c)$.
\textbf{(V2)} By the chain rule: $\sigma(p) = f'(g(x)) \cdot \frac{d}{dx}[g(x)] = f'(g(x)) \cdot \sigma(c)$. Thus $\sigma(c) \in \text{factors}(\sigma(p))$.
\textbf{(V3)} By Definition~\ref{def:expr_tree}, $T_p = \text{node}(f, [T_c])$, so $T_c$ is a direct subtree.
\end{proof}

\end{document}
